\chapter{Belajar Tanpa Label}
\label{ch:unsup}

Pada musim panas 2025 saya mengikuti dua kuliah yang saling melengkapi. Kuliah unsupervised
learning dari Sepp Hochreiter memakai naskah panjang yang bergerak dari PCA dan ICA sampai
biclustering, metode yang lahir dari analisis data ekspresi gen. Kuliah machine learning and
pattern classification memberi sisi praktisnya: pengklasifikasi dasar, seleksi fitur, dan cara
membangun sistem pengenal pola dari awal sampai akhir.

Sebagian besar data di dunia tidak berlabel. Rekaman sensor, foto di ponsel, dan hasil
pengukuran ekspresi ribuan gen datang tanpa keterangan tentang apa isinya. Unsupervised learning
mencari struktur di dalam data seperti itu: arah-arah yang penting, sumber-sumber tersembunyi,
kelompok-kelompok alami, atau peta berdimensi rendah yang mempertahankan tetangga. Hasilnya sering
menjadi langkah pertama sebelum model supervised dibangun.

\section{Dua cara melihat PCA}

PCA sudah muncul di \cref{ch:cv} sebagai eigenfaces. Kuliah unsupervised menekankan bahwa ada dua
cara setara untuk menurunkannya. Cara pertama mencari arah $\bm v$ dengan varians proyeksi
terbesar, $\max_{\|\bm v\|=1}\bm v^\top\bm\Sigma\bm v$, yang jawabannya vektor eigen kovarians
dengan nilai eigen terbesar. Cara kedua mencari subruang berdimensi $q$ yang meminimalkan galat
rekonstruksi kuadrat rata-rata ketika data diproyeksikan ke subruang itu lalu dikembalikan. Kedua
masalah punya jawaban yang sama, karena varians total tetap, dan varians yang tidak tertangkap oleh
proyeksi adalah tepat galat rekonstruksinya.

Penurunan cara pertama memakai pengali Lagrange dari \cref{ch:optimasi}. Kita memaksimalkan $\bm v^\top\bm\Sigma\bm v$
dengan kendala $\bm v^\top\bm v=1$. Menyamakan gradien fungsi tujuan dengan $\lambda$ kali gradien kendala memberi
$2\bm\Sigma\bm v=2\lambda\bm v$, yaitu $\bm\Sigma\bm v=\lambda\bm v$: setiap titik stasioner adalah vektor eigen. Varians
proyeksi pada vektor eigen itu $\bm v^\top\bm\Sigma\bm v=\lambda$, sehingga yang terbesar adalah vektor eigen dengan
nilai eigen terbesar.

Contoh dua dimensi membuatnya konkret. Untuk data dengan kovarians
$\bm\Sigma=\begin{pmatrix}2&1\\1&2\end{pmatrix}$, nilai eigennya $3$ dan $1$, dengan vektor eigen $(1,1)/\sqrt2$ dan
$(1,-1)/\sqrt2$. Data membentang paling lebar sepanjang diagonal, dan komponen utama pertama menangkap $3/(3+1)=75$
persen varians. Memproyeksikan ke satu dimensi itu membuang seperempat varians, dan galat rekonstruksi rata-ratanya
tepat sama dengan nilai eigen yang dibuang, yaitu satu.

PCA unik hanya sampai tanda setiap komponen, dan kalau dua nilai eigen sama, arah-arahnya bisa
diputar bebas di dalam subruang mereka. Komponen utama tidak berkorelasi satu sama lain, tetapi
tidak berarti independen. PCA juga linear. Untuk struktur yang melengkung, \istilah{kernel PCA}
\citep{scholkopf1998} memetakan data secara implisit ke ruang fitur berdimensi tinggi lewat sebuah
kernel, lalu menjalankan PCA di sana. Yang dibutuhkan hanya matriks kernel antar titik data, bukan
pemetaannya secara eksplisit.

\section{Memisahkan sumber: ICA}

Di sebuah pesta, dua orang berbicara bersamaan, dan dua mikrofon di sudut ruangan merekam
campuran suara keduanya dengan proporsi yang berbeda. Bisakah kita mendapatkan kembali suara
masing-masing orang tanpa tahu proporsi campurannya? Inilah \istilah{cocktail party problem}, dan
\istilah{independent component analysis} \citep{hyvarinen2000} menjawabnya. Modelnya
$\vx=\mA\bm s$, dengan $\bm s$ sumber yang saling independen dan $\mA$ matriks campuran yang tidak
diketahui. Kita mencari matriks pemisah $\mW$ sehingga $\bm y=\mW\vx$ sedekat mungkin dengan
independen.

Kuncinya adalah teorema limit pusat. Campuran dari beberapa sumber independen cenderung lebih
mirip Gauss daripada sumber-sumber aslinya. Jadi kita bisa membalik logikanya: cari arah proyeksi
yang hasilnya paling \emph{tidak} Gauss. Ketidak-Gaussan diukur dengan kurtosis atau negentropi. Kurtosis berlebih, $\E[y^4]/\E[y^2]^2-3$,
bernilai nol untuk distribusi Gauss, negatif untuk distribusi yang ``rata'' seperti distribusi seragam ($-1{,}2$), dan
positif untuk distribusi yang runcing dengan ekor tebal seperti distribusi Laplace ($+3$). Suara manusia cenderung
runcing: kebanyakan waktu hampir diam, sesekali keras. Campuran dua suara sudah lebih dekat ke Gauss, dan kurtosisnya
lebih kecil. Memaksimalkan nilai mutlak kurtosis berarti mencari arah yang paling mirip satu suara saja.
Langkah praktisnya dimulai dengan \istilah{whitening}, yaitu mentransformasi data agar
kovariansnya menjadi matriks identitas. Setelah itu, yang tersisa hanya mencari rotasi, dan
algoritme seperti FastICA atau Infomax menemukannya. ICA tidak bisa memulihkan urutan dan skala
sumber, dan ia gagal kalau lebih dari satu sumber berdistribusi Gauss, karena campuran Gauss
yang diputar tetap Gauss.

Perbedaan dengan PCA mudah dilihat pada sinyal suara. PCA menemukan arah varians terbesar, yang
biasanya masih campuran kedua pembicara. ICA menemukan arah yang membuat setiap keluaran terdengar
seperti satu orang.

\section{Analisis faktor}

Analisis faktor mirip probabilistic PCA (\cref{ch:mlat}), $\vx=\mW\vz+\bm\mu+\bm\varepsilon$, dengan
satu perbedaan penting: noise $\bm\varepsilon$ punya kovarians diagonal $\bm\Psi$ yang boleh berbeda
untuk setiap variabel. Setiap variabel teramati punya ``noise pribadinya'' sendiri, dan faktor
laten $\vz$ menjelaskan korelasi di antara mereka. Model ini tidak punya solusi tertutup, sehingga
dilatih dengan EM. Dalam data biologis, di mana setiap gen atau setiap probe punya tingkat noise
yang berbeda, asumsi ini jauh lebih masuk akal daripada noise yang sama untuk semua.

\section{Peta berdimensi rendah}

Kadang yang dicari bukan sumber, melainkan peta: gambaran dua dimensi dari data berdimensi tinggi
yang mempertahankan sebanyak mungkin hubungan kedekatan. Kuliah membahas sederet metode dengan
filosofi yang berbeda.

\istilah{Multidimensional scaling} mencari titik-titik di bidang yang jarak-jaraknya sedekat mungkin
dengan jarak asli. \istilah{Non-negative matrix factorization} menguraikan data menjadi bagian-bagian
yang hanya boleh dijumlahkan, tidak boleh dikurangkan, sehingga komponennya sering bisa ditafsirkan
sebagai bagian-bagian nyata, misalnya bagian wajah atau topik dalam teks. \istilah{Locally linear
embedding} \citep{roweis2000} menganggap setiap titik bisa direkonstruksi secara linear dari
tetangganya, lalu mencari koordinat rendah dimensi yang mempertahankan bobot rekonstruksi itu.
Isomap \citep{tenenbaum2000} mengganti jarak lurus dengan jarak di sepanjang graf tetangga, sebuah
pendekatan jarak geodesik, sehingga data yang tergulung seperti gulungan kue bisa dibentangkan.
\istilah{Self-organizing map} \citep{kohonen1982} menempatkan grid neuron di ruang data dan menariknya
ke arah data sambil menjaga neuron yang bertetangga di grid tetap berdekatan. Metode yang paling
populer untuk visualisasi saat ini adalah t-SNE dan UMAP \citep{mcinnes2018}, yang mencocokkan
peluang ketetanggaan di ruang asli dengan di ruang peta.

Peta seperti ini menyenangkan dilihat, tetapi mudah disalahbaca. Jarak antar kelompok di peta
t-SNE tidak bermakna kuantitatif, dan ukuran kelompok bisa diperbesar atau diperkecil oleh
algoritmenya. Kuliah AI and Visualization kembali ke masalah ini di \cref{ch:xai}.

\section{Pengelompokan}

\subsection{k-means}

Algoritme $k$-means \citep{lloyd1982} bergantian antara dua langkah: tetapkan setiap titik ke pusat
terdekat, lalu pindahkan setiap pusat ke rata-rata titik-titik anggotanya. Setiap langkah tidak
menaikkan jumlah kuadrat jarak titik ke pusatnya, sehingga algoritme pasti berhenti, meski belum
tentu di solusi terbaik. Ambil empat titik di garis bilangan, 1, 2, 10, dan 11, dengan pusat awal 1
dan 2. Pada langkah pertama, titik 1 masuk ke pusat pertama, sedangkan 2, 10, dan 11 masuk ke pusat
kedua, yang lalu pindah ke rata-ratanya, $23/3\approx7{,}67$. Pada langkah kedua, titik 2 sekarang
lebih dekat ke pusat pertama, sehingga kelompoknya menjadi $\{1,2\}$ dan $\{10,11\}$, dengan pusat
$1{,}5$ dan $10{,}5$. Langkah berikutnya tidak mengubah apa pun. Kalau pusat awalnya buruk, hasil
akhirnya bisa buruk juga, sehingga $k$-means biasanya dijalankan beberapa kali dari awal acak yang
berbeda.

\subsection{Model campuran dan EM}

$k$-means adalah versi keras dari model campuran Gauss. Pada model campuran, setiap titik punya
peluang menjadi anggota setiap kelompok, dan kelompoknya boleh lonjong dengan kovarians sendiri.
Algoritme EM bergantian antara menghitung peluang keanggotaan (langkah E) dan memperbarui rata-rata,
kovarians, dan bobot setiap komponen dengan rata-rata berbobot (langkah M):
\begin{equation}
\gamma_{nk}=\frac{\pi_k\,\N(\vx_n\mid\bm\mu_k,\bm\Sigma_k)}{\sum_j\pi_j\,\N(\vx_n\mid\bm\mu_j,\bm\Sigma_j)},\qquad
\bm\mu_k=\frac{\sum_n\gamma_{nk}\vx_n}{\sum_n\gamma_{nk}},\qquad \pi_k=\frac1N\sum_n\gamma_{nk}.
\end{equation}
Langkah E adalah teorema Bayes: peluang titik $n$ berasal dari komponen $k$ sebanding dengan bobot komponen itu dikali
seberapa mungkin titik itu dihasilkannya. Langkah M adalah estimasi maksimum likelihood biasa, hanya dengan setiap titik
dihitung sebagian untuk setiap komponen. Setiap putaran EM tidak pernah menurunkan likelihood, alasan yang sama dengan
batas bawah ELBO di \cref{ch:mlat}. Ketika kovarians semua
komponen dibuat sama dan kecil sekali, keanggotaan menjadi nol atau satu, dan EM berubah menjadi
$k$-means.

Berapa banyak kelompok yang sebaiknya dipakai? Jumlah kuadrat jarak ke pusat selalu turun kalau $k$ ditambah, sampai
nol ketika setiap titik menjadi kelompoknya sendiri, sehingga ukuran itu tidak bisa dipakai langsung. Metode siku mencari
titik di mana penurunannya tiba-tiba melandai. Koefisien silhouette membandingkan, untuk setiap titik, jarak rata-rata
ke anggota kelompoknya sendiri, $a$, dengan jarak rata-rata ke kelompok terdekat lainnya, $b$, sebagai $(b-a)/\max(a,b)$.
Nilainya mendekati satu kalau titik itu jelas berada di kelompok yang benar. Untuk contoh empat titik tadi, titik 1 punya
$a=1$ dan $b=(9+10)/2=9{,}5$, sehingga silhouette-nya $8{,}5/9{,}5\approx0{,}89$. Untuk model campuran, kriteria informasi
seperti BIC menimbang likelihood terhadap jumlah parameter.

\subsection{Hierarki dan kemiripan}

Pengelompokan hierarkis tidak memilih jumlah kelompok di awal. Versi aglomeratif mulai dengan setiap
titik sebagai kelompoknya sendiri, lalu berulang kali menggabungkan dua kelompok terdekat. Arti
``terdekat'' menentukan bentuk hasilnya. \istilah{Single linkage} memakai pasangan titik terdekat dan
cenderung menghasilkan rantai panjang, \istilah{complete linkage} memakai pasangan terjauh dan
menghasilkan kelompok yang padat, dan \istilah{average linkage} berada di antaranya. Hasilnya adalah
dendrogram yang bisa dipotong di ketinggian mana pun. Metode berbasis kemiripan, termasuk
pengelompokan spektral, memandang data sebagai graf kemiripan dan memotongnya di tempat ikatan
paling lemah, sebuah ide yang sudah kita temui pada segmentasi gambar.

\subsection{Biclustering}

Data ekspresi gen berbentuk matriks: baris adalah gen, kolom adalah sampel jaringan atau kondisi.
Sekelompok gen sering hanya bekerja sama di sebagian sampel, misalnya hanya di jaringan tumor
tertentu. Mengelompokkan seluruh baris atau seluruh kolom akan melewatkan pola ini.
\istilah{Biclustering} mencari submatriks, yaitu sekelompok baris yang berperilaku mirip pada
sekelompok kolom. FABIA \citep{hochreiter2010}, yang dikembangkan di Linz, merumuskannya sebagai
model faktor dengan prior yang jarang sekali (\istilah{sparse}), sehingga setiap faktor hanya melibatkan
sedikit gen dan sedikit sampel. Inilah salah satu titik pertemuan pertama antara machine learning
dan ilmu hayati di buku ini, dan kita akan kembali ke data seperti ini di \cref{ch:hayati}.

\section{Dari fitur ke pengklasifikasi}

Kuliah pattern classification membahas jalur lengkap sebuah sistem pengenal pola: ambil sinyal
mentah, ubah menjadi fitur, pilih fitur yang berguna, latih pengklasifikasi, lalu evaluasi dengan
jujur. Contoh yang dipakai banyak berasal dari audio, seperti mendeteksi kejadian akustik dalam
rekaman.

Pengklasifikasi dasar yang dibahas punya karakter yang berbeda-beda. $k$-nearest neighbour tidak
belajar apa pun dan hanya mengingat data latih, lalu meminjam label tetangga terdekat. Naive Bayes
menganggap semua fitur independen bersyarat terhadap kelas, asumsi yang hampir selalu salah tetapi
sering bekerja dengan baik. Analisis diskriminan linear \citep{fisher1936} mencari arah proyeksi
yang memisahkan rata-rata kelas sejauh mungkin relatif terhadap sebaran di dalam kelas. Untuk dua kelas,
memaksimalkan rasio $(\bm w^\top(\bm\mu_1-\bm\mu_2))^2/(\bm w^\top\bm S_W\bm w)$ memberi arah
$\bm w\propto\bm S_W^{-1}(\bm\mu_1-\bm\mu_2)$, dengan $\bm S_W$ matriks sebaran di dalam kelas. Kalau sebaran itu bulat,
arahnya sekadar garis yang menghubungkan kedua rata-rata. Kalau sebaran lonjong, arahnya dimiringkan menjauhi arah
sebaran yang lebar, karena perbedaan rata-rata di arah itu mudah tertutup noise. Pohon
keputusan memotong ruang fitur dengan pertanyaan ya atau tidak yang berurutan. Pertanyaan di setiap simpul dipilih
yang paling mengurangi ketidakmurnian label, misalnya entropi $H=-\sum_cp_c\log_2p_c$. Simpul dengan delapan contoh
positif dan delapan negatif punya entropi satu bit. Kalau sebuah pertanyaan membaginya menjadi $(7,1)$ dan $(1,7)$,
entropi setiap anak sekitar $0{,}54$ bit, sehingga perolehan informasinya sekitar $0{,}46$ bit. Pertanyaan yang membagi
menjadi $(4,4)$ dan $(4,4)$ tidak memberi informasi sama sekali. Pohon tunggal mudah overfitting, dan random forest
\citep{breiman2001} merata-ratakan banyak pohon yang dilatih pada sampel dan fitur acak. Support
vector machine \citep{cortes1995} mencari bidang pemisah dengan margin terlebar, dan dengan kernel ia
bisa memisahkan kelas yang tidak terpisah secara linear. Ide marginnya adalah regularisasi dalam
bentuk geometris: di antara semua bidang yang memisahkan data latih, pilih yang paling jauh dari
titik mana pun. Jarak sebuah titik ke bidang $\bm w^\top\vx+b=0$ adalah $|\bm w^\top\vx+b|/\|\bm w\|$. Kalau skala $\bm w$
diatur sehingga titik-titik terdekat memenuhi $y_n(\bm w^\top\vx_n+b)=1$, lebar margin di kedua sisi menjadi
$2/\|\bm w\|$. Memaksimalkan margin sama dengan meminimalkan $\tfrac12\|\bm w\|^2$ dengan kendala
$y_n(\bm w^\top\vx_n+b)\ge1$ untuk semua $n$, masalah kuadrat dengan kendala yang diselesaikan lewat bentuk dualnya
(\cref{ch:optimasi}). Di bentuk dual, data hanya muncul sebagai hasil kali dalam $\vx_n^\top\vx_m$, dan mengganti hasil
kali itu dengan kernel memberi SVM nonlinear tanpa pernah menghitung pemetaan fiturnya.

Seleksi fitur bisa dilakukan dengan tiga cara. Metode filter menilai setiap fitur sendiri-sendiri,
misalnya dari korelasinya dengan label. Metode wrapper mencoba berbagai himpunan fitur dengan
melatih model sungguhan. Metode embedded memilih fitur sebagai bagian dari pelatihan, seperti lasso
yang mendorong bobot fitur yang tidak berguna ke nol. Kuliah juga memperkenalkan jaringan saraf
dengan cara pandang praktis: kapan lapisan dense cukup, kapan konvolusi lebih cocok, dan bagaimana
transfer learning serta pretraining self-supervised mengurangi kebutuhan data berlabel.

\section{Di luar kelas}

Metode tanpa label muncul di hampir setiap bidang yang disentuh buku ini. Dalam ilmu hayati,
pengelompokan profil ekspresi gen memisahkan subtipe kanker yang tidak terlihat di bawah
mikroskop, dan biclustering menemukan modul gen yang aktif bersama. Dalam pemrosesan sinyal, ICA
memisahkan sinyal jantung ibu dan janin dari elektroda yang sama, dan PCA memampatkan data sensor
sebelum dikirim. Dalam simulasi, PCA dengan nama proper orthogonal decomposition mereduksi medan
aliran berjuta sel menjadi beberapa puluh mode (\cref{ch:penemuan}). Dalam sistem rekomendasi,
faktorisasi matriks yang mirip analisis faktor menemukan selera tersembunyi pengguna
(\cref{ch:rekomendasi}). Di semua tempat itu, pertanyaannya sama: struktur apa yang tersembunyi di
balik angka-angka ini?

\section{Perkembangan riset}

Batas antara unsupervised dan supervised learning kini kabur karena \istilah{self-supervised learning}: label dibuat dari
data itu sendiri. SimCLR \citep{chen2020simclr} melatih encoder agar dua versi augmentasi dari gambar yang sama berdekatan
dan versi dari gambar lain berjauhan, loss kontrastif yang sama dengan CLIP di \cref{ch:hayati}. Studi sistematisnya
menemukan bahwa susunan augmentasi, satu proyeksi nonlinear tambahan, dan batch yang besar adalah kunci keberhasilannya.
Masked autoencoder \citep{he2022mae} mengambil jalan lain: sembunyikan tiga perempat potongan gambar secara acak, lalu
latih jaringan merekonstruksi piksel yang hilang. Karena encoder hanya melihat potongan yang terlihat, pelatihannya juga
tiga kali lebih cepat.

I-JEPA \citep{assran2023} berargumen bahwa merekonstruksi piksel membuang kapasitas untuk detail yang tidak penting.
Ia meramalkan \emph{representasi} dari blok-blok yang disembunyikan, bukan pikselnya, dari satu blok konteks, sehingga
yang dipelajari lebih semantik. DINOv2 \citep{oquab2023} menunjukkan bahwa dengan data yang dikurasi dengan cermat dan
model miliaran parameter, fitur self-supervised bisa dipakai langsung untuk banyak tugas tanpa penyetelan ulang. Analisis
faktor dan PCA di bab ini mencari struktur laten dengan model linear. Metode-metode ini mencari struktur yang sama dengan
jaringan dalam, dan menjadi titik awal model fondasi di bidang visi, citra medis, dan sinyal sensor.

\sumberbab{Bab ini mengikuti naskah kuliah Machine Learning: Unsupervised Techniques (PCA dan kernel
PCA, ICA, analisis faktor, metode proyeksi, clustering, biclustering) dan kuliah Machine Learning and
Pattern Classification (pengklasifikasi dasar, seleksi dan konstruksi fitur, jaringan saraf dalam
praktik). Bacaan pendamping: \citet{bishop2006}, \citet{hastie2009}, dan \citet{duda2001}.}
